
%(BEGIN_QUESTION)
% Copyright 2012, Tony R. Kuphaldt, released under the Creative Commons Attribution License (v 1.0)
% This means you may do almost anything with this work of mine, so long as you give me proper credit

Les og lag en disposisjon av underkapitlet ``Current Transformers'' i delen ``Electrical Sensors'' i kapitlet ``Electric Power Measurement og Control'' i læreboken {\it Lessons In Industrial Instrumentation}. Noter sidetall der viktige illustrasjoner, fotografier, ligninger, tabeller og andre relevante detaljer finnes. Forbered deg på en faglig diskusjon med lærer og medstudenter om begrepene og eksemplene i lesingen.

\underbar{file i01244}
%(END_QUESTION)




%(BEGIN_ANSWER)


%(END_ANSWER)





%(BEGIN_NOTES)

En strømtransformator trinntransformerer og isolerer høy linjestrøm til et nivå som er trygt for panelinstrumenter som amperemetre. Fast CT-forhold gir en proporsjonal representasjon av linjestrøm. CT-er er ofte toroidale, der én gjennomføring av linjelederen utgjør én primærvinding.

5 A er industristandard for å representere 100\% linjestrøm. Denne standardiseringen gjør at ulike CT-fabrikat og panelinstrumenter er kompatible. Noen CT-er bruker likevel lavere sekundærstrøm (f.eks. 1 A).

\vskip 10pt

Siden en CT oppfører seg som en strømkilde når linjestrøm flyter, må sekundærviklingen aldri åpnes. Da kan svært høy spenning oppstå, fordi CT-ens strømnedtrapping tilsvarer spenningsopptrapping. Derfor er sikker metode ved frakobling av spenningssatt CT å {\it kortslutte} sekundærsiden. Strømkilder skades ikke av kortslutning slik spenningskilder gjør.

\vskip 10pt

Effektbrytere og annet høyspenningsutstyr har ofte CT-er montert ved terminalgjennomføringer (``bushings''), for praktisk måling av strøm inn og ut av utstyret. Noen svært høyspente CT-er monteres i enden av lange isolatorer for avstand til jord.







\vskip 20pt \vbox{\hrule \hbox{\strut \vrule{} {\bf Forslag til sokratisk diskusjon} \vrule} \hrule}

\begin{itemize}
\item{} Hva er hensikten med å bruke CT i elektrisk effektinstrumentering?
\item{} Hva er standard utgangssignalområde for en kraftlinje-CT?
\item{} Se orienteringen av sekundærviklingene i en ``donut''-CT. Forklar hvorfor viklingene må legges i denne retningen, og ikke parallelt med kjernens omkrets.
\item{} Forklar i detalj faren ved åpne sekundærviklinger på strømtransformatorer under drift.
\item{} Forklar hvordan omsetningsforholdet kan endres i en vindus-CT.
\item{} Anta at effektlederen føres tre ganger gjennom åpningen i en CT merket 400:5. Beregn effektivt omsetningsforhold.
\item{} Anta at effektlederen føres fire ganger gjennom åpningen i en CT merket 600:5. Beregn effektivt omsetningsforhold.
\item{} Har potensialtransformatorer (PT) samme fare ved åpen sekundærvikling som strømtransformatorer (CT)? Hvorfor eller hvorfor ikke?
\item{} Ubrukte CT-sekundærviklinger bør kortsluttes. Forklar hvorfor dette ikke skader en spenningssatt CT, og hvorfor det faktisk er den sikreste løsningen.
\end{itemize}










\vfil \eject

\noindent
{\bf Forberedende quiz:}

Hvorfor skal du aldri åpne sekundærviklingen på en strømtransformator?

\begin{itemize}
\item{} Det kan bli vanskelig å etablere forbindelsen igjen
\vskip 5pt 
\item{} Det kan oppstå farlig høy spenning i den åpne kretsen
\vskip 5pt 
\item{} Transformatorens målenøyaktighet vil få spennfeil
\vskip 5pt 
\item{} Transformatorens målenøyaktighet vil få nullpunktsfeil
\vskip 5pt 
\item{} Det kan føre til at overstrømsvernet feilaktig registrerer høy strøm
\vskip 5pt 
\item{} Det anses rett og slett som dårlig praksis
\end{itemize}



%INDEX% Leseoppgave: Lessons In Industrial Instrumentation, strømtransformatorer

%(END_NOTES)
